% Research appendix fragment. Requires amsmath, amssymb, amsthm.
% New alternative presentation; no claim that its checker is smaller.
\section{A progressive binder presentation of iPRA}
\label{app:progressive-binder}

This section gives a different, complete presentation of iPRA in which
every binder introduction has the same progress form.  The induction
restriction moves to elimination.  Boolean case analysis must be
retained explicitly; deleting it is not justified here.  The result is a
mathematical alternative, not a measured improvement over the main
checker.

\subsection{Numerical syntax and complete rule inventory}

Numerical terms are variables, $0$, $St$, and
$R(n,a;u,v.h)$.  Only the fixed finite update body $h$ binds $u,v$;
other parameters are held fixed.  The numerical computation schemas are
\[
 R(0,a;u,v.h)=a,\qquad
 R(Sn,a;u,v.h)=h[n/u,R(n,a;u,v.h)/v].
 \tag{PR}
\]
All substitutions avoid capture, and the successor substitution is
simultaneous.  Equivalently use a staged first-order symbol for each
fixed body.  This is ordinary counter-aware primitive recursion, so
zero, successor, projections, composition, and arbitrary fixed PR
descriptions compile directly.  There is no uniform evaluator over
quantified descriptions.

Object formulas have grammar
\[
 A::=(s=t)\mid\mathsf H_{x,y}A.
\]
Internal clauses are $\mathsf P_yA(y)$, where $A$ is an object formula.
Contexts are reusable finite lists of objects and clauses.  An object
or clause is available when it is a context member; a proved object is
also available.  Clauses are not object goals.  Put $\bot_*=(0=S0)$.

Besides the computation schemas, the complete primitive rule inventory is
\[
 \frac{\Gamma\mathrel{\triangleright}s=t\quad
       \Gamma\vdash A(s)}{\Gamma\vdash A(t)}
 \quad\mathsf{Tr},
\]
\[
 \frac{\Gamma,B(t,z)\vdash B(t,Sz)}
      {\Gamma\vdash\mathsf H_{x,y}B(x,y)}
 \quad\mathsf{HI},
\]
\[
 \frac{\Gamma\mathrel{\triangleright}\mathsf H_{x,y}B(x,y)
       \quad\Gamma,\mathsf P_yB(a,y)\vdash A}
      {\Gamma\vdash A}
 \quad\mathsf{HE},
\]
\[
 \frac{\Gamma\mathrel{\triangleright}\mathsf P_yB(y)
       \quad\Gamma\vdash B(t)}
      {\Gamma\vdash B(St)}
 \quad\mathsf{Step},
\]
\[
 \frac{\Gamma\mathrel{\triangleright}\mathsf P_yE(y)
       \quad\Gamma\vdash E(0)}
      {\Gamma\vdash E(t)}
 \quad\mathsf{Iter},
 \qquad E\text{ a literal equation},
\]
and Boolean cases
\[
 \frac{\Gamma,\operatorname{sg}(t)=0\vdash A\quad
       \Gamma,\operatorname{sg}(t)=1\vdash A}
      {\Gamma\vdash A}
 \quad\mathsf{BoolCases}.
\]
Here $1=S0$ and $\operatorname{sg}(n)=R(n,0;u,v.1)$, so
$\operatorname{sg}(0)=0$ and $\operatorname{sg}(Sn)=1$.
In HI, $z$ is fresh for $\Gamma,t$ and the passive parameters.
In HE, $a$ is fresh for $\Gamma,A$ and the passive parameters.
BoolCases has no eigenvariable.  The displayed local
assumptions and clauses are discharged.  There is no primitive
hypothesis-conclusion rule: context membership supplies availability
only, and object identity is recovered by the expansion below.
There are no other logical axioms or induction rules.

The intended ordinary meanings are
\[
 (\mathsf P_yB)^\circ=\forall y(B^\circ(y)\to B^\circ(Sy)),
\]
\[
 (\mathsf H_{x,y}B)^\circ
   =\exists x\,\forall y\bigl(B^\circ(x,y)\to B^\circ(x,Sy)\bigr).
 \tag{Meaning}
\]
HI is unrestricted logical introduction.  Iter is the only use of
induction, and it is explicitly restricted to literal equations.

\subsection{Elementary derived operations}

Structural weakening, contraction, exchange, substitution, renaming,
and assumption cut are admissible by induction on derivations, with
freshening of eigenparameters before substitution.  The equation
$R(0,t;u,v.0)=t$ and transport derive $t=t$, followed by symmetry,
transitivity, and numerical congruence.

An equation assumption $s=t$ transports reflexivity $s=s$ to $s=t$.
For a paired assumption
$\mathsf H B$, HE supplies a fresh witness $a$ and $\mathsf P B(a,-)$.
Introduce the same paired object at $a$; the induction hypothesis on
formula depth recovers its assumed matrix $B(a,z)$, and Step yields
$B(a,Sz)$.  Thus every object identity has a finite expansion.

Literal-equation induction is derivable.  Given $E(0)$ and a derivation
of $E(Sz)$ from $E(z)$, HI with an unused witness gives
$\mathsf H_{x,y}E(y)$.  HE supplies its progress clause, and Iter
gives $E(t)$.  Conversely Iter is sound in ordinary iPRA by equation
induction under the retained assumption $E(0)$.

\paragraph{Boolean cases recover constructor cases.}
Put $P(n)=R(n,0;u,v.u)$ and $C(n,a,b)=R(n,a;u,v.b)$.
Their equations give $P0=0$, $P(Sn)=n$ and the usual selector
endpoints.  Literal-equation induction proves
\[
 C(\operatorname{sg}(x),0,x)=x,\qquad
 C(\operatorname{sg}(x),x,S(Px))=x.
 \tag{Boolean coverage identities}
\]
Each base is a direct zero computation, and each successor step is a
direct successor computation that does not use its induction hypothesis.
Consequently $\operatorname{sg}(t)=0$ gives $t=0$ by the first
identity, and $\operatorname{sg}(t)=1$ gives $t=S(Pt)$ by the second.
It follows that the constructor rule
\[
 \frac{\Gamma,t=0\vdash A\quad\Gamma,t=Sz\vdash A}
      {\Gamma\vdash A}
 \quad\mathsf{Cases}
\]
is admissible, with $z$ fresh for $\Gamma,t,A$: use BoolCases, and
in its second branch substitute $P(t)$ for the fresh $z$ in the
supplied successor derivation.  Conversely constructor Cases proves
BoolCases by the sign equations.  Thus the two case interfaces have
exactly the same consequences.  The missing principle is isolated to
the two-value range of sign; equation induction already supplies the
needed numerical identities.  In the rest of the proof, Cases denotes
this derived constructor rule.

\paragraph{Binder congruence.}
Suppose $A(x,y)$ and $B(x,y)$ are uniformly interderivable with
$x,y$ absent from the retained context.  From $\mathsf H A$, HE gives
the progress clause of $A(a,y)$.  To introduce $\mathsf H B$ at
the same witness, convert the local $B(a,z)$ to $A(a,z)$, apply Step,
and convert $A(a,Sz)$ to $B(a,Sz)$.  Reverse the conversions for the
opposite direction.  Hence $\mathsf H$ preserves uniform equivalence.
This argument needs neither Cases nor compound induction.

\subsection{Gates and selectors on arbitrary formulas}

Let $C(n,a,b)$ be the fixed PR numerical selector with
$C(0,a,b)=a$ and $C(Sn,a,b)=b$.  It is directly represented by
$R(n,a;u,v.b)$.  Define a numerical gate on formulas by
\[
 G_v(s=t):=(C(v,0,s)=C(v,0,t)),\qquad
 G_v(\mathsf H_{x,y}A):=\mathsf H_{x,y}G_vA,
\]
freshening bound variables first.  External formula induction, numerical
computation, and binder congruence give
\[
 \vdash G_0A,\qquad G_{Sv}A\dashv\vdash A.
\]
For the first assertion at a binder, HI ignores its local assumption
and uses the induction hypothesis for the successor matrix.

Padding is again available:
\[
 A\dashv\vdash \mathsf H_{u,y}G_yA
 \qquad(u,y\text{ fresh for }A).
 \tag{Padding-H}
\]
Forwards HI ignores its local assumption and uses the positive-gate
equivalence.  Backwards HE supplies a progress clause; $G_0A$ is a
theorem, so one Step gives $G_1A$ and hence $A$.

The selector on formulas is defined exactly by finite structural
recursion:
\[
 M_v(s=t,r=w):=(C(v,s,r)=C(v,t,w)),
\]
\[
\begin{aligned}
 M_v(\mathsf H_{x,y}A,\mathsf H_{x,y}B)
      &:=\mathsf H_{x,y}M_v(A,B),\\
 M_v(E,\mathsf H_{x,y}B)&:=\mathsf H_{x,y}M_v(G_yE,B),\\
 M_v(\mathsf H_{x,y}A,E)&:=\mathsf H_{x,y}M_v(A,G_yE),
\end{aligned}
\]
where $E$ is an equation.  Bindings are aligned by fresh renaming.
The maximum binder depth decreases in every recursive call.  Computation,
padding, and congruence give
\[
 M_0(A,B)\dashv\vdash A,\qquad M_{Sv}(A,B)\dashv\vdash B.
 \tag{Selector-H}
\]
Cases therefore permits choosing either endpoint at an arbitrary control,
including when the retained context mentions that control.  In particular
proofs of both $A$ and $B$ give a proof of $M_t(A,B)$.

\subsection{All intuitionistic logical connectives}

Use $\mathsf T=(0=0)$.  Choose fixed PR functions $p,q,d$ with
\[
 p(0)=0,\quad p(Sn)=C(p(n),1,0),
\]
\[
 q(0)=0,\quad q(Sn)=C(p(n),q(n),S(q(n))),
\]
\[
 d(0)=0,\quad d(Sn)=SS(d(n)).
\]
They are parity, integer half, and doubling in the standard model.
The following equations are proved internally by literal-equation
induction and their defining computations:
\[
 p(d(n))=0,\quad p(Sd(n))=1,\quad
 q(d(n))=n,\quad q(Sd(n))=n.
 \tag{Clock}
\]
First induct on the single equation $p(d(n))=0$.  At a successor,
two parity updates send zero to one and back to zero.  The identity
$p(Sd(n))=1$ then follows by one computation.  Next induct on the
single equation $q(d(n))=n$.  The two half updates at $d(n)$ and
$Sd(n)$ first leave the value unchanged and then increment it, using
the already proved parity equations.  Finally $q(Sd(n))=n$ follows
by its defining zero-parity update.  Thus the argument uses two
single-equation inductions in sequence, not conjunction induction or
induction on any compound formula.

Define the logical macros
\[
\begin{aligned}
 A\Rightarrow_H B
   &:=\mathsf H_{u,z}M_z(A,B),\\
 \exists_H x\,A(x)
   &:=\mathsf H_{x,z}G_zA(x),\\
 A\wedge_H B
   &:=\mathsf H_{u,z}G_zM_{p(z)}(A,B),\\
 A\vee_H B
   &:=\exists_H x\,M_x(A,B),\\
 \forall_H x\,A(x)
   &:=\mathsf H_{u,z}M_{p(z)}(\mathsf T,A(q(z))),\\
 \bot_H&:=(0=S0).
\end{aligned}
 \tag{H-macros}
\]
All displayed unused coordinates and binders are fresh.  The universal
macro uses alternating true positions: its matrix is true at $d(n)$
and equivalent to $A(n)$ at $Sd(n)$.
These macros are effective translations of fixed finite syntax, not
numerical truth predicates or internal quantification over formula codes.

\paragraph{Implication.}
Given $\Gamma,A\vdash B$, apply HI to the implication macro.
Its local assumption is $M_z(A,B)$.  Cases splits $z=0$ and
$z=Sw$.  The first branch provides $A$ and then $B$; the second
already provides $B$.  In both branches the successor matrix is $B$.
For elimination, HE supplies the progress clause; $A$ is its zero
matrix, and Step at zero gives its first positive matrix $B$.

\paragraph{Existential quantification.}
From $A(t)$ use HI at witness $t$, ignoring the local gate assumption.
For elimination, HE gives a fresh witness $a$ and its progress clause.
The zero gate is a theorem; Step at zero gives $A(a)$.  The ordinary
fresh-witness elimination convention follows directly from HE.

\paragraph{Conjunction.}
If both $A$ and $B$ are proved, Cases and the selector endpoints prove
every successor matrix, so HI applies.  Conversely HE gives the
progress clause.  Its zero gate is a theorem, so Step at zero gives
the matrix at $1$, namely $B$.  Step once more gives the matrix at $2$,
namely $A$, using $p(1)=1$ and $p(2)=0$.

\paragraph{Disjunction.}
The two introductions choose existential witnesses zero and one.
For elimination, use existential elimination at a fresh witness $a$;
Cases splits its zero and successor alternatives.  The selector endpoints
then give the respective supplied branch assumptions.  The final
conclusion is arbitrary, with the usual retained context.

\paragraph{Universal quantification.}
Suppose $\Gamma\vdash A(x)$ with $x$ fresh for $\Gamma$.
After numerical substitution, both $\mathsf T$ and $A(q(Sz))$ are
proved.  Cases and the selector endpoints therefore prove the successor
matrix of the universal macro; HI may ignore its local assumption.
For specialization, HE provides the progress clause.  The matrix at
$d(t)$ is a theorem by Clock and its true endpoint.  Step at $d(t)$
gives the matrix at $Sd(t)$, which Clock identifies with $A(t)$.
This proof does not iterate a compound progress clause.

\paragraph{Falsity and coverage.}
Under $0=1$, numerical computation and transport give
$s=C(0,s,t)=C(1,s,t)=t$ for all numerical $s,t$.
External formula induction then gives every object: for an H node,
HI ignores the local assumption and uses the induction hypothesis on
its successor matrix.  Thus explosion is derived.  Cases also proves
the ordinary coverage macro
$(n=0)\vee_H\exists_H v\,(n=Sv)$ by its two evident introductions.

Consequently every ordinary intuitionistic natural-deduction rule is
derived, including all eigenvariable restrictions.  No excluded middle
for arbitrary quantified formulas has been used.

\subsection{Whole-domain bi-interpretability}

For a transparent comparison, let $T_H$ be ordinary intuitionistic
first-order logic with equality in the staged PR signature, with the
displayed PR computation schemas, literal-equation induction,
constructor coverage, and $(0=1)\to\bot$.  The native rules translate
to $T_H$ under Meaning: HI, HE, and Step are ordinary logical rules;
Iter is equation induction; BoolCases is the sign range consequence
of coverage.  Conversely the native
system derives all the ordinary logical rules, coverage, and equation
induction, with separation translating to the identity implication
$(0=1)\Rightarrow_H(0=1)$.

The usual actual-equality characteristic construction for
PR computation plus coverage, separation, and literal-equation
induction proves every quantifier-free induction instance.  In
particular this is the same arithmetic completion theorem proved in
the arithmetic appendix of the main report.  Its hypotheses have all
been established here.  Thus $T_H$ and iPRA have the same consequences
in their shared staged PR language; the domain is the entire numerical
domain and equality is actual equality.

It remains to check the native formula composite, not merely the
logical macros.  For every fixed native matrix $B$ one has
\[
 \mathsf H_{x,y}B(x,y)\dashv\vdash
 \exists_H x\,\forall_H y\,
       \bigl(B(x,y)\Rightarrow_H B(x,Sy)\bigr).
 \tag{H-composite}
\]
Forwards, HE supplies a fresh $a$ and its progress clause.  At fresh
$y$, assume $B(a,y)$ and use Step to obtain $B(a,Sy)$.  Introduce
implication and universal macros, then introduce the existential at
$a$, and close HE.  Backwards, eliminate the existential at fresh
$a$.  Specialize its universal at the fresh HI eigenparameter $z$.
Under HI's assumption $B(a,z)$, implication elimination gives
$B(a,Sz)$.  HI at witness $a$ and existential elimination finish.
The analogous clause composite is the universal progress formula.

In ordinary $T_H$, the gate, selector, and Clock equations prove the
intended meanings of every macro.  The universal macro is worth
emphasizing: its even-to-odd progress implications give every instance
$A(n)$, while all instances of $A$ make every progress implication
true.  No induction on $A$ appears in this equivalence.

External induction on formula syntax, using binder congruence and
H-composite, proves both interpretation composites.  Numerical staged
descriptions are identified with themselves; if one passes to a
different conventional PR enumeration, literal-equation induction
uniquely identifies the corresponding fixed recursions.  Both
internally definable domain isomorphisms have graph $x=y$.
The progressive-binder calculus is therefore equiconsistent and
whole-domain bi-interpretable with full first-order iPRA.

\paragraph{Limit of the reduction.}
This presentation has an unrestricted single HI rule, but it pays for
that uniformity with explicit BoolCases and a distinction between Step and
literal-equation Iter.  No theorem here makes Cases redundant, and no
smaller complete Python checker for this presentation has been
established.  It is therefore an alternative logical basis, not evidence
that the measured frontier has been beaten.
